\section{Tensions et enjeux \'epist\'emologiques}\label{sec:tensions}

Les deux positions exposées à la section précédente partagent une limite commune: elles traitent la question des IAG comme un problème de régulation sans interroger la validité des critères d'évaluation sur lesquels elles s'appuient. 
C'est précisément ce que nous allons faire ici. 
Les IAG ne créent pas les tensions suivantes: elles les révèlent, parce qu'elles poussent à leur point de rupture des présupposés que le système scolaire n'avait jamais eu besoin d'expliciter.


\subsection*{Tension 1: l'appartenance cognitive d'un travail produit avec IAG}

La première tension porte sur l'\emph{appartenance cognitive}: à qui appartient un travail produit avec l'aide d'une IAG?\@
La question semble simple, mais elle révèle une ambiguïté au cœur de la définition même de la compétence.%
~\cite{tardif2006evaluation} définit la compétence comme la mobilisation de ressources \emph{internes et externes}: dans cette perspective, l'IAG est une ressource externe légitime, au même titre qu'un dictionnaire, un logiciel spécialisé ou une base de données. 
Refuser ce statut à l'IAG exige de justifier pourquoi certaines ressources externes sont légitimes et d'autres non --- justification que le système scolaire n'a pas encore fournie de façon cohérente.

Mais la compétence suppose aussi une \emph{intériorisation}: on ne peut pas dire qu'un étudiant est compétent en argumentation s'il est incapable de construire un raisonnement sans IAG, même s'il sait très bien lui soumettre un prompt. 
Le critère de l'intériorisation n'est pas arbitraire: il renvoie à l'idée que la compétence doit être \emph{transférable} à des situations nouvelles, y compris des situations où l'outil est absent.
C'est précisément ce que la définition de~\cite{perrenoud2011construire} appelle la capacité à agir face à une \og{}famille de situations\fg{}: si cette capacité dépend entièrement de la présence de l'outil, elle n'est pas généralisable.

La tension est d'autant plus aiguë que les IAG sont \emph{invisibles dans leurs traces}: contrairement à un logiciel de statistiques dont on voit le code, ou à une calculatrice dont les résultats intermédiaires trahissent l'usage, le recours à une IAG ne laisse aucune marque formelle dans le texte produit.
L'évaluation a toujours présumé une cohérence entre performance affichée et capacités mobilisées: les IAG brisent cette présomption. 
La réponse à cette tension n'est pas d'interdire la ressource, mais de concevoir des évaluations où le processus est observable --- l'oral, la mise en situation, le commentaire réflexif sur sa propre démarche --- plutôt que de s'appuyer sur le seul rendu final.

\subsection*{Tension 2: l'évaluation mesure-t-elle encore une compétence individuelle?}

La deuxième tension touche au fondement \emph{individualiste} de l'évaluation scolaire. 
Notre système certifie des individus: chaque réussite est attribuée à une personne, en tant que trace de ses capacités propres. 
Ce présupposé est cohérent avec l'idéologie méritocratique qui légitime le système scolaire --- l'idée que la hiérarchie des résultats reflète une hiérarchie des efforts et des aptitudes.

\cite{pluet1971pierre} ont montré que ce présupposé était déjà partiellement fictif: la réussite scolaire co-dépend du capital culturel familial, c'est-à-dire d'une ressource inégalement distribuée que l'école fait passer pour un mérite individuel. % chktex 2
Les IAG radicalisent cette déconstruction d'une façon nouvelle. 
Ce que Bourdieu analysait comme la transmission tacite d'un style, d'une façon de penser et d'écrire, peut désormais être partiellement simulé par un outil accessible à tous. 
La question n'est donc pas seulement: \og{}est-ce que cet étudiant a triché?\fg{} mais: \og{}qu'est-ce que nous cherchions réellement à mesurer en demandant ce type de travail?\fg{}

La réponse qui s'impose est un déplacement de l'unité d'évaluation: plutôt que de certifier un produit fini dont on ne peut plus garantir l'authenticité individuelle, évaluer \emph{la façon dont l'apprenant interagit avec les outils à sa disposition}. 
La compétence à l'ère des IAG n'est peut-être pas la compétence sans outil, mais la compétence \emph{dans} un écosystème outillé --- ce qui implique de rendre cet écosystème visible dans le dispositif d'évaluation lui-même.

\subsection*{Tension 3: redistribution de l'autorité du savoir}

La troisième tension concerne l'\emph{autorité épistémique}. 
Dans le système scolaire traditionnel, l'enseignant est le détenteur et le médiateur du savoir légitime; l'évaluation mesure la conformité de la production de l'apprenant aux critères institutionnels. 
Cette architecture suppose que l'institution scolaire est le lieu où se valide ce qui compte comme savoir.

Les IAG perturbent cette architecture en introduisant une autorité épistémique concurrente, non institutionnelle. 
Les apprenants les interrogent pour vérifier leurs réponses, structurer leurs idées, valider leurs hypothèses: elles exercent de fait une autorité épistémique sans être inscrites dans les dispositifs institutionnels de validation. 
Ce phénomène n'est pas entièrement nouveau --- Internet et Wikipédia avaient déjà initié cette redistribution --- mais les IAG l'intensifient qualitativement: là où un moteur de recherche renvoyait à des sources que l'apprenant devait évaluer, une IAG produit directement une réponse synthétisée, avec une apparence d'autorité et de cohérence.

\cite{lepage2023recension} montrent que la transformation du rôle de l'enseignant est un fil conducteur de toute l'histoire de l'IA en éducation, mais que les IA génératives l'accélèrent considérablement. % chktex 2
Le rôle de transmetteur du savoir ne disparaît pas, mais il se déplace vers celui de \emph{formateur à l'évaluation critique des sources et des productions}: l'enseignant devient la personne qui apprend aux apprenants à questionner ce que l'IAG leur dit, à identifier ses angles morts, à distinguer une réponse plausible d'une réponse vérifiée. 
C'est une transformation profonde, qui exige de repenser la formation des enseignants autant que les dispositifs d'évaluation.

\subsection*{Tension 4: la crise des critères}

La quatrième tension, la plus profonde, est une \emph{crise des critères}. 
Les exercices évalués --- dissertation, mémoire, rapport --- étaient justifiés par l'idée qu'ils mobilisaient des compétences complexes et authentiquement humaines. 
Si une IAG peut produire un texte formellement satisfaisant selon ces critères, c'est que les critères eux-mêmes sont insuffisants pour distinguer une compétence authentiquement développée d'une simulation compétente.

Cette crise n'est pas entièrement nouvelle: des didacticiens ont depuis longtemps signalé le risque du \og{}contrat didactique perverti\fg{}~\citep{leroux2015evaluer}, où l'apprenant apprend à satisfaire les critères formels de l'évaluation sans développer les compétences visées.
Ce comportement est d'ailleurs structurellement analogue à ce que fait une IAG pendant son propre entraînement: optimiser une fonction d'objectif définie par des évaluateurs humains, sans que cet objectif capture nécessairement ce qu'on cherche vraiment à produire. 
Les IAG portent la logique du contrat didactique perverti à son point de rupture: elles rendent la simulation d'une production compétente techniquement accessible à tous.

Mais cette rupture est aussi une opportunité épistémologique. 
Si les IAG peuvent satisfaire nos critères actuels, c'est une invitation à formuler de \emph{meilleurs critères} --- des critères qui capturent ce que les outils ne peuvent pas faire à notre place. 
Plusieurs candidats se dégagent de l'analyse: la réflexivité sur son propre processus d'apprentissage, la capacité à évaluer et corriger une production générée, le jugement situé dans un contexte singulier, l'engagement éthique avec le savoir. 
Ces dimensions ne sont pas nouvelles dans la réflexion pédagogique; les IAG leur donnent simplement une urgence qu'elles n'avaient pas encore.